\section{ISDB-Tb sobre cable}
\label{sec:isdbtb_cable}

La propuesta es distribuir por una red de cable canales modulados en ISDB-Tb, la misma norma de la TDA, en lugar de canales analógicos PAL-N o digitales DVB-C. Cada canal de $6 \, [MHz]$ de la grilla de cable lleva una señal OFDM idéntica a la de una emisora terrestre, con uno o varios programas. Se usa sobre redes coaxiales o HFC de $75 \, [\Omega]$, tanto en cableoperadores como en instalaciones cerradas de edificios, hoteles o clínicas, y también puede entregarse por RF Overlay en una red FTTH (\autoref{sec:rf_overlay}).

La ventaja principal es que los televisores vendidos en Argentina en los últimos años traen un sintonizador ISDB-Tb incorporado, porque es la norma de la TDA. Al recibir los canales en esa norma, el televisor los sintoniza y los decodifica directamente, sin decodificador externo, y el operador puede reemplazar los canales analógicos por canales digitales en alta definición sin instalar un decodificador por televisor. Además, en un canal de $6 \, [MHz]$ entran varios programas digitales donde antes había uno analógico, lo que libera espectro para DOCSIS.

\subsection{Equipos}
En la cabecera, los receptores toman las señales de origen (satélite, fibra o IP), los codificadores H.264 comprimen las que llegan sin comprimir o en otro formato, y un multiplexor-modulador ISDB-Tb arma los flujos de transporte y genera las portadoras OFDM en los canales elegidos; un combinador de RF las suma al resto de la grilla. La red es la misma red de cable, con sus amplificadores y derivadores, en la que hay que respetar los niveles y la planicidad de la respuesta, porque una inclinación del espectro dentro del canal deja algunas subportadoras con menor relación señal a ruido. En el abonado alcanza con un televisor con sintonizador ISDB-Tb conectado al coaxial.

\subsection{Compatibilidad con los receptores}
Como la señal cumple la norma terrestre, los fabricantes declaran compatibilidad con cualquier receptor ISDB-Tb. VideoSwitch indica que su equipo genera las portadoras según ARIB STD-B31 y ABNT NBR 15601 y que el uso de ISDB-Tb en redes cableadas ``permite la transmisión de televisión digital HD sin necesidad de agregar receptores externos al TV'' \cite{videoswitch_dmm}. Del análisis propio, no de la documentación del fabricante, surgen dos condiciones: el televisor debe poder sintonizar las frecuencias que use el operador, ya que algunos receptores podrían explorar sólo los canales UHF 21 a 51 de la TDA, y como los televisores comunes no tienen módulo de acceso condicional, los programas deben distribuirse sin cifrar.

\subsection{Ejemplo de solución comercial: VideoSwitch DMM-3100}
VideoSwitch, empresa argentina de Buenos Aires, fabrica el modulador-multiplexor DMM-3100, ``ISDB-Tb Over Cable'', destinado a cableoperadores e instalaciones cerradas como hoteles y sanatorios (\autoref{tab:dmm3100}). Según la prensa del sector, la empresa desarrolló esta tecnología para ofrecer HD y liberar ancho de banda en la red de cable, y su primera venta fue a Telecom \cite{prensario_videoswitch}. Otros fabricantes ofrecen soluciones equivalentes, como el módulo Luminato ISDB-T de Teleste \cite{teleste_isdbt}.

\begin{table}[H]
    \centering
    \caption{Características del VideoSwitch DMM-3100 según la hoja de datos del fabricante \cite{videoswitch_dmm}.}
    \label{tab:dmm3100}
    \begin{tabular}{p{4cm} p{10.2cm}}
        \toprule
        \textbf{Característica} & \textbf{Valor declarado} \\
        \midrule
        Modulación & COFDM según ARIB STD-B31 y ABNT NBR 15601 \\
        Portadoras y salidas & Hasta 8 portadoras de $6 \, [MHz]$, en 2 salidas de $75 \, [\Omega]$ con hasta 4 portadoras adyacentes cada una, ágiles de $50$ a $863 \, [MHz]$ \\
        Potencia de salida & De $-18$ a $+3 \, [dBm]$ \\
        Entradas & 4 ASI y 1 GbE (UDP/IP unicast o multicast) \\
        Multiplexación & Hasta 32 servicios por portadora, con generación automática de tablas PSI/SI \\
        Formato & 1 unidad de rack de 19", fuentes redundantes, monitoreo por SNMP \\
        \bottomrule
    \end{tabular}
\end{table}

No se encontró una norma específica de ISDB-Tb sobre cable: la propuesta aplica la norma terrestre al medio cableado.

\subsection{Comparación con DVB-C}
En un canal de $6 \, [MHz]$, ISDB-Tb llega como máximo a $23{,}2 \, [Mbps]$, contra $38{,}5 \, [Mbps]$ de DVB-C con 256-QAM (\autoref{sec:dvbc}), porque las herramientas de robustez de OFDM (intervalo de guarda, pilotos y codificación convolucional) consumen capacidad y no son necesarias en el cable. A cambio, ISDB-Tb funciona con el parque de televisores existente, mientras que DVB-C requiere decodificadores salvo en los televisores que lo incluyen. Para un cableoperador argentino que quiere digitalizar los canales básicos sin entregar decodificadores, la compatibilidad pesa más que la eficiencia espectral.
